\documentclass[11pt,a4paper]{ltjsarticle}
\usepackage[margin=25mm]{geometry}
\usepackage{amsmath,amssymb,amsthm,mathtools,booktabs,longtable,array}
\usepackage[haranoaji]{luatexja-preset}
\usepackage[unicode,hidelinks]{hyperref}
\usepackage{xurl}
\newtheorem{theorem}{定理}[section]
\newtheorem{proposition}[theorem]{命題}
\newtheorem{lemma}[theorem]{補題}
\newtheorem{corollary}[theorem]{系}
\theoremstyle{definition}
\newtheorem{definition}[theorem]{定義}
\newtheorem{remark}[theorem]{注意}
\newcommand{\ZZ}{\mathbb Z}
\newcommand{\GG}{\mathbb G}
\newcommand{\Gr}{\operatorname{Gr}}
\newcommand{\GL}{\operatorname{GL}}
\newcommand{\Ext}{\operatorname{Ext}}
\newcommand{\Tor}{\operatorname{Tor}}
\newcommand{\Hom}{\operatorname{Hom}}
\newcommand{\gr}{\operatorname{gr}}
\newcommand{\Per}{\operatorname{Per}}
\newcommand{\rk}{\operatorname{rank}}
\newcommand{\NF}{\operatorname{NF}}
\newcommand{\qgr}{\operatorname{qgr}}
\title{非可換重み付き射影空間\texorpdfstring{$\mathbb P(1,1,2,3)$}{P(1,1,2,3)}の分類に向けて\\
\large 整域部分のOre構造と一般対象の七窓表示}
\author{}
\date{2026年11月19日\quad 現行研究資料 v12\\
\small 完全分類および新規性の最終確認は未完成}
\begin{document}
\maketitle
\noindent 本稿は未完成の一般AS正則$\mathbb Z$-代数分類の現行研究資料であり、公表用統合論文は\path{papers/nc-p1123.tex}である。以下のgraded部分の記述は統合論文を参照する。
本資料の対象は$n=2$、重み$(1,1,2,3)$に固定する。
任意$n$の整域への拡張では、単一のgraded代数への帰着は$n=1$の最小周期2の例により否定される。
その完全な証明は統合論文\texttt{genidx:thm:period-two-counterexample}にある。
この例は添字付きOre表示を持つため、任意$n$の添字付き二段Ore表示の網羅性とは区別する。
$n=1$の全対象の添字付きOre表示・裸の同型・周期は、統合論文
\texttt{genidx:thm:one-light-ore}、\texttt{genidx:thm:one-light-classification}で決定した。
一般の添字付きOre網羅性で残る範囲は$n\geq3$である。
\begin{abstract}
最小生成元の重みが$(1,1,2,3)$で、関係が語長二次かつ重み付き斉次である
四次元AS正則$\ZZ$-代数の添字保持同型による分類を調べる。
広い対象$\mathcal Q$では、AS双対性と多項式増大から語長Koszul性、AS添字7、
語長を保つ7周期同型が従う。所期のHilbert級数を持つことと標準解消を持つことは同値である。
また整域性を仮定せず$r_2(i)=r_5(i)=1$を強制でき、
純jump 1部分は各行の次元が$n+1$となる。
有限幅の数値的分岐は、混合関係の制約と積共変ベクトルによって排除される。
混合関係数$r_3(i)=r_4(i)=2$も行列Hilbert級数の成長と890列の厳密な有限整数検査で強制でき、
全$\mathcal Q$で標準解消と標準Hilbert級数が従う。後半は計算補助の証明である。

さらに$\ZZ$-整域である対象については、標準Hilbert級数も1周期性も仮定する必要がない。
整域性による関係数の上界と七行Euler式から標準型を強制し、
重み2の正規生成元族が与える2周期性と7周期性を組み合わせて1周期性を導く。
対応するgraded代数は量子平面またはJordan平面からの2段Ore拡大であり、
左右noetherian性も結論として得られる。この構造定理の完全証明は
公表用統合論文\path{papers/nc-p1123.tex}に置く。
さらに、選んだ語長二次表示の重み2・語長1の生成元族$z_i$が全て左右正則であることは、
$A$が$\ZZ$-整域であることと同値である。7周期により、調べる生成元は七つに絞れるが、
各乗法写像の単射性は全次数で必要である。
最小周期7または非定数の関係射影階数が現れるとすれば、
少なくとも一つの$z_i$に零因子がある部分に限られる。

整域部分の二つのOre表示が裸の添字保持同型で一致するための必要十分条件は、
一つの可逆な$2\times2$行列と非零スカラーに関する有限系へ帰着した。
対角量子多項式族では、三係数トーラスの明示的な対合による商が、
全特殊値を含む裸の同型類の集合を与える。
本稿はこれらの結果の要点に加え、一般対象の七窓表示、traceの選択、
二つの基準族と同型不変量を保持する。
全対象の整域性、零因子を持つ部分の存在と分類、
一般Oreデータの軌道の標準形およびモジュライ幾何は未解決である。
構造定理と既存文献の完全な照合および新規性の最終確認も未完了である。
\end{abstract}
\tableofcontents

\section{問題の位置づけと先行研究}
本稿では$k$を標数零の代数閉体とする。
重み付き多項式環$k[x,y,z,t]$に
$\deg(x,y,z,t)=(1,1,2,3)$を与えたものを可換基準とする。
語長を別に数えると、$tx$と$z^2$はともに語長2、重み4である。
この二つの次数を分けることが、通常の四次元二次AS正則代数との比較の出発点となる。

Auroux--Katzarkov--Orlov\cite[\S2, Proposition 2.6]{AKO}は、任意の重みの
対角量子多項式環を扱い、そのAS正則性とGorenstein parameterが重みの和となることを証明した。
同論文Theorem 2.12の例外生成列は、今回の対角族では
$\mathcal O,\ldots,\mathcal O(6)$の7項となる。
また、可換重み付き環の$\qgr$は重み付き射影\emph{スタック}の連接層の圏に対応する
\cite[Proposition 2.3, Corollary 2.4]{AKO}。
本稿はこれらを既知の基準族とし、全対象が対角族で尽くされるとは仮定しない。

$\ZZ$-代数の一般的な設定と1周期性はVan den Bergh\cite[\S3--4]{VdB}に従う。
graded代数$B$に対応する$\check B$を考えると、Sierra\cite[Theorem 1.2]{Sierra}により
Zhang twistは既に$\check B$の同型として同一視される。
したがってgraded同型類、指定周期構造付きの同型類、裸の添字保持同型類、
$\qgr$同値を分ける必要がある。

AS双対性からExt圏がFrobeniusとなる一般原理は
Mart\'inez-Villa--Solberg\cite[Lemma 6.1(b), Theorem 6.2]{MVS}にある。
本稿では必要な有限射影解消の議論を重み付き添字のまま記述し、
語長二次性と多項式増大によってAS添字を7に絞り、7移動と有限表示を導く。
一般的なKoszul--Frobenius原理そのものを新規結果とはしない。
Ore拡大の正則性にはLiu--Wang--Wu\cite[Theorem 3.3]{LWW}の一般理論も使えるが、
本稿の具体的族では二次双対とtraceを明示し、仮定を直接確認する。

同じHilbert級数を持つ代数が同じ分類対象であるとは限らない。
Lu--Palmieri--Wu--Zhang\cite{LPWZ}の二つの次数1生成元を持つ四次元正則代数では、
関係次数が3と4で、解消のBetti型は$(1,2,2,2,1)$となる。
そのHilbert級数は本稿の標準型と一致するが、最小生成元と語長二次性が異なる。
このため本稿ではHilbert級数だけで対象を定義しない。

整域部分の構造定理は、標準Hilbert級数・1周期性・noetherian性を結論として導く点で、
それらを仮定したgraded部分のOre定理より適用範囲が広い。
一般のFrobenius Ext圏、1周期からgraded代数への復元、正則正規元の商は既知の道具であり、
本研究では七行Euler式による関係数の強制と正規元族の2周期性を結び付ける。
これらの完全証明は統合論文に置き、本研究資料は七窓による一般対象の分類に必要な補足を保持する。
これらの個別結果の完全な文献上の新規性確認、および全分類は未完成である。

\section{分類対象と二つの次数}
正連結で局所有限な$\ZZ$-代数$A$は
\[
A_{ii}=k,\qquad A_{ij}=0\quad(i>j),\qquad
A_{ij}A_{j\ell}\subset A_{i\ell}
\]
を満たす。$I=A_+=\bigoplus_{i<j}A_{ij}$、$V=I/I^2$と置く。
添字差$j-i$を重みと呼ぶ。$P_i=e_iA$、$S_i=P_i/P_iI$、
$L_i=Ae_i/IAe_i$を、それぞれ右射影・右単純・左単純加群とする。

\begin{definition}[広い対象]\label{def:target}
$\mathcal Q$は次を満たす$A$の添字保持同型類の集合とする。
\begin{enumerate}
\item 各$i$で$\dim V_{i,i+1}=2$, $\dim V_{i,i+2}=\dim V_{i,i+3}=1$、
それ以外の$V_{ij}$は零である。
\item 添字を保つ語長斉次二次表示
$A\simeq T_{\ZZ}(V)/(R)$, $R\subset V\otimes_{\ZZ}V$を持つ。
ここで$V\to I$の持上げは表示の選択であり、分類に固定しない。
\item 左右のglobal dimensionは4、各単純加群は有限生成射影による有限解消を持ち、
行・列の累積次元は多項式的に増大する。
\item ある添字の全単射$\ell:\ZZ\to\ZZ$について
$\Ext_A^p(S_i,A)=0$ ($p\ne4$),
$\Ext_A^4(S_i,A)\simeq L_{\ell(i)}$であり、左側にも対応するAS双対性がある。
\end{enumerate}
比較のため、各行Hilbert級数が
$((1-t)^2(1-t^2)(1-t^3))^{-1}$となる部分を$\mathcal H\subseteq\mathcal Q$とする。
\end{definition}
添字を一括して移動した対象で割ることはしない。
またdomain、noetherian、PBW、CY、1周期性を定義に追加しない。
語長斉次表示を選ぶと$A$は語長$n$と重み$d$の二変数次数を持つ。
重み付き斉次であっても語長の混じる関係は、本稿のquadratic条件には含めない。

\begin{definition}[標準重み付きKoszul型]\label{def:standard}
$\mathcal S$は$\mathcal Q$のうち、語長に関してKoszulであり、
ASの添字移動が$\ell(i)=i+7$で、各右単純加群に次の最小解消があるものとする。
\begin{equation}\label{eq:resolution}
\begin{gathered}
0\longrightarrow P_{i+7}\longrightarrow
P_{i+4}\oplus P_{i+5}\oplus P_{i+6}^{\oplus2}\longrightarrow
P_{i+2}\oplus P_{i+3}^{\oplus2}\oplus P_{i+4}^{\oplus2}\oplus P_{i+5}\\
\longrightarrow P_{i+1}^{\oplus2}\oplus P_{i+2}\oplus P_{i+3}
\longrightarrow P_i\longrightarrow S_i\longrightarrow0.
\end{gathered}
\end{equation}
ホモロジー次数$p$の項の生成元は語長$p$に置く。
左側にも対応する解消とAS双対性を要求する。
\end{definition}

\begin{proposition}[数値型]\label{prop:hilbert}
$A\in\mathcal S$の各行の二変数Hilbert級数は
\[
H_i(u,t)=\frac{1}{(1-ut)^2(1-ut^2)(1-ut^3)}.
\]
従って$\mathcal S\subseteq\mathcal H$である。
\end{proposition}
\begin{proof}
解消\eqref{eq:resolution}のEuler和の係数は、$w=(1,1,2,3)$について
\[
\prod_{r=0}^3(1-ut^{w_r})
=1-u(2t+t^2+t^3)+u^2(t^2+2t^3+2t^4+t^5)
-u^3(t^4+t^5+2t^6)+u^4t^7
\]
となる。まだ行が同じであると仮定してはいけない。
Euler等式は$H_i$を正の重みだけ先に進んだ行$H_{i+d}$の低い係数から決める。
定数項は全て1であるため、重みに関する帰納法により全行の係数は一致し、
上の多項式を逆数に持つ。$u=1$とすれば主張を得る。
\end{proof}
特に
\[
H(t)^{-1}=1-2t+t^3+t^4-2t^6+t^7.
\]
この一変数式では生成元と関係等の寄与が相殺される。
これを逆に読んで最小解消を一意に復元することはできない。
後の系\ref{cor:hilbert-standard}で$\mathcal S=\mathcal H$を示す。
定理\ref{thm:all-standard-summary}により$\mathcal H=\mathcal Q$も成立する。
整域部分では、主論文に計算非依存の別証明もある。

\section{同型の線型化}
\begin{proposition}\label{prop:linearization}
$A=T_{\ZZ}(V)/(R)$と$B=T_{\ZZ}(W)/(R')$を、正の重みを持つ語長斉次二次表示とする。
$A\simeq B$が添字保持同型として成立することと、各$V_{ij}\to W_{ij}$の線型同型が存在し、
誘導される二次テンソル写像が$R$を$R'$へ移すことは同値である。
\end{proposition}
\begin{proof}
添字保持同型は$A_+$とその全ての冪を保つ。
語長斉次表示では$I^n=\bigoplus_{m\geq n}A^{(m)}$なので、選んだ表示から
$\gr_I A\simeq A$を得る。各添字成分では語長は重み以下であり、この等式に完備化は不要である。
同型に付随次数を取れば語長を保つ同型を得て、その語長1部分が所要の線型写像となる。
語長2の積の核を取れば関係も対応する。
逆はテンソル代数の普遍性と、関係イデアルが二次関係で生成されることによる。
\end{proof}

\begin{remark}
この命題は同型の\emph{存在}についての主張であり、全ての同型が線型とは言っていない。
例えば可換例では$z\mapsto z+x^2$が重みを保つ自己同型である。
従って線型群作用の商群oidを、非線型自己同型も持つ裸の代数のmoduli stackと同一視しない。
\end{remark}

\section{AS添字7の強制と七周期}
\subsection{AS双対性・Koszul性・添字の強制}
この部分の完全証明は\path{papers/nc-p1123.tex}の
\texttt{idx:lem:duality-period}と\texttt{idx:thm:index-seven}に統合した。
ここでは七窓の議論で用いる声明・記法と証明の要点を記す。
前向きの規約で$E_{ij}^p=\Ext_A^p(S_i,S_j)$とする。

\begin{lemma}\label{lem:frobenius}
$A\in\mathcal Q$についてYoneda積
\[
E_{ij}^p\otimes E_{j,\ell(i)}^{4-p}\longrightarrow E_{i,\ell(i)}^4\simeq k
\]
は完全対合である。
\end{lemma}
\begin{proof}[証明の参照と要点]
有限生成射影解消$P_\bullet\to S_i$の$A$-双対は、AS条件により
$L_{\ell(i)}$の左射影解消を反転した複体である。
$\Hom_A(P_\bullet,M)\simeq M\otimes_A\Hom_A(P_\bullet,A)$から
$D\Ext_A^p(S_i,M)\simeq\Ext_A^{4-p}(M,S_{\ell(i)})$を得る。
接続準同型との自然性により、この対合はYoneda合成に一致する。
詳細は統合論文\texttt{idx:lem:duality-period}、一般原理は
\cite[Lemma 6.1(b), Theorem 6.2]{MVS}を参照する。
\end{proof}

\begin{proposition}[語長Koszul性と一定のAS添字]\label{prop:general-koszul}
$A\in\mathcal Q$は語長Koszulであり、ある正整数$L$について$\ell(i)=i+L$となる。
さらに語長を保つ同型$A\simeq A(L)$が存在する。
ここで$A(s)_{ij}=A_{i+s,j+s}$である。
\end{proposition}
\begin{proof}[証明の参照と要点]
$E^1,E^2$の語長は1,2である。$E^2$同士の完全対合からtopの語長は4、
次いで$E^1,E^3$の完全対合から$E^3$の語長は3と分かる。
従って最小解消は語長線型である。
trace $\lambda_i:E_{i,\ell(i)}^4\to k$を選ぶと
\begin{equation}\label{eq:nakayama}
\lambda_i(ab)=\lambda_j\bigl(bN_{ij}(a)\bigr),\qquad
N_{ij}:E_{ij}^p\xrightarrow{\sim}E_{\ell(i),\ell(j)}^p
\end{equation}
によってNakayama同型が定まる。
$E^1$の各行で次元2の唯一の成分はjump 1なので$\ell(i+1)=\ell(i)+1$である。
Yoneda積の結合性から$N(ac)=N(a)N(c)$となり、
Koszul双対$E=A^!$の語長1部分の逆双対によって$A\simeq A(L)$を得る。
完全証明は統合論文\texttt{idx:lem:duality-period}にある。
\end{proof}

\begin{lemma}[二次元因子の区間]\label{lem:segment}
各$U_0,\ldots,U_{n-1}$を二次元とし、隣接位置ごとの関係空間
$D_j\subset U_j\otimes U_{j+1}$の次元を高々1とする。
全テンソル積を関係挿入空間の和で割った空間の次元は2以上である。
\end{lemma}
\begin{proof}
最初の因子で独立な二つの共変ベクトルを選び、それぞれを順次延長する。
次の非零共変ベクトルに課される条件は高々一本なので延長できる。
得られる二つの分解可能汎関数は全関係を消し、最初の因子の独立性から独立である。
\end{proof}

\begin{theorem}[AS添字の強制]\label{thm:general-index}
定義\ref{def:target}の全対象についてAS添字は$L=7$である。
\end{theorem}
\begin{proof}[証明の参照と要点]
$r_d(i)=\dim R_{i,i+d}$と置く。完全対合から
\begin{equation}\label{eq:relation-symmetry}
r_d(i)=r_{L-d}(i+d)
\end{equation}
が成立する。$E^1,E^3$の重みから$6\leq L\leq10$である。
$L\geq9$では$r_2=0$となり、純jump 1の語が自由に増大する。
$L=6$では$r_5=r_6=0$となり、jump 1生成元と$zz$を消した商に指数増大が残る。
$L=8$では、$www$のsyzygyを除くAS対合から
$r_2(j)+r_2(j+3)\leq1$となる。
零関係の辺で区切った純jump 1部分に補題\ref{lem:segment}を適用すると指数増大する。
いずれも多項式増大に反する。各除外と隠れた関係が生じないことの完全証明は
統合論文\texttt{idx:thm:index-seven}にある。
\end{proof}

\begin{corollary}\label{cor:hilbert-standard}
$\mathcal H=\mathcal S$である。すなわち$\mathcal Q$の対象が所期の重み付きHilbert級数を持つことと、
標準重み付き最小解消を持つことは同値である。
\end{corollary}
\begin{proof}
$\mathcal S\subseteq\mathcal H$は命題\ref{prop:hilbert}で示した。
逆に$A\in\mathcal H$ではKoszul性とAS添字7が既知なので、Euler等式から
\[
\sum_d r_d(i)t^d
=H(t)^{-1}-1+(2t+t^2+t^3)+(t^4+t^5+2t^6)-t^7
=t^2+2t^3+2t^4+t^5.
\]
左右のAS対合により左解消も標準型となる。
\end{proof}

\subsection{純jump 1部分と外側の関係数}
\begin{proposition}\label{prop:degree-one-dichotomy}
$A\in\mathcal Q$に対し、純jump 1生成部分を$B$、
$a_i=\dim R_{i,i+2}$と置く。次のどちらか一方が成立する。
\begin{enumerate}
\item 全ての$i$で$a_i=1$であり、$\dim B_{i,i+n}=n+1$である。
\item 全ての$i$と$n\geq14$について$B_{i,i+n}=0$である。
\end{enumerate}
これは整域性を仮定しない全対象についての中間的な二分岐である。
下の定理\ref{thm:outer-relation-ranks}で第二の枝を排除する。
\end{proposition}
\begin{proof}[証明の所在と要点]
完全証明は\path{papers/nc-p1123.tex}の
\texttt{idx:prop:degree-one-dichotomy}にある。
bar複体の「重み$-$語長」が零の部分を取ると$B$のbar複体となる。
AS添字7での$\operatorname{Tor}_3,\operatorname{Tor}_4$の重みから、
$B$はKoszulで、左右の単純加群の射影次元は高々2である。
従って$b_{i,n}=\dim B_{i,i+n}$は
\[
b_{i,n}=2b_{i,n-1}-a_{i+n-2}b_{i,n-2},
\qquad b_{i,0}=1,\quad b_{i,1}=2
\]
を満たす。$a_i\in\{0,1,2\}$は7周期である。
$a_{i-1}\ne0$を選ぶと、実際の二つの係数列
$b_{i,n}$と$b_{i-1,n+1}$の初期ベクトルは独立であるため、
七段の移行行列の全固有値は絶対値1以下となる。
全ての$a_i$が正なら行列式から全て1となり、零を含む場合は
七段が奇数であることからtraceの偶数性を使って移行行列の二乗が零となる。
これが上の二分岐を与える。
有限個の数値表の検査によって証明したものではない。
\end{proof}
\begin{theorem}[外側の関係数の強制]\label{thm:outer-relation-ranks}
全ての$A\in\mathcal Q$と$i\in\ZZ$について
\[
r_2(i)=r_5(i)=1,\qquad \dim B_{i,i+n}=n+1
\]
である。
\end{theorem}
完全証明は統合論文\texttt{idx:thm:outer-relation-ranks}にある。
有限幅の枝を仮定すると、混合関係の三重重なりとモノミアル商の成長から、
一周期中に$a_i=2$の位置がちょうど一つあり、$a_i=0$も現れる。
零関係の辺でtensor鎖を分割すると、各ブロックの二次元関係は高々一つである。
代数閉性からそのannihilatorに非零の積共変ベクトルが存在し、
一次元関係を消すように左右へ延長できるため、任意長の$B_{i,i+n}$が非零となる。
これが有限幅と矛盾する。数値条件だけでは排除できない具体的な係数表と、
その方法上の障害は\path{research/notes_ja.tex}に置く。

関係空間の次元$r_2(i)=1$と、その生成tensorの行列階数2は異なる主張である。
後者、$B$の整域性およびAS正則次元2は、全$\mathcal Q$ではまだ証明していない。
例えば一般の二次代数$k\langle x,y\rangle/(xy)$はHilbert級数$(1-u)^{-2}$を持ち、
Koszulで単純加群の射影次元は2だが零因子を持つ。
これは$\mathcal Q$の構成例ではなく、純jump 1部分の数値からの推論を制限する比較例である。
$r_3(i)=2$の強制は次の定理で得るが、$z_i$の全次数での左右正則性は未解決である。

\begin{theorem}[全対象の標準数値型：計算補助の証明]
\label{thm:all-standard-summary}
全ての$A\in\mathcal Q$について$r_3(i)=r_4(i)=2$であり、
$\mathcal Q=\mathcal H=\mathcal S$である。
特に全ての行の二変数Hilbert級数は
$((1-ut)^2(1-ut^2)(1-ut^3))^{-1}$となる。
\end{theorem}
証明は主論文\texttt{idx:thm:mixed-relation-ranks}と\texttt{idx:cor:all-standard}にある。
七行の二変数Hilbert級数の正係数と多項式増大から行列式恒等式を導き、
$\delta_i=r_3(i)-2\in\{-2,-1,0,1,2\}$、$\sum\delta_i=0$、
$\delta_i+\delta_{i+1}\leq1$という必要条件を得る。
残る890列の整数行列式を網羅計算すると、零になるのは定数零列だけである。
\path{computations/verify_mixed_ranks.py}と\path{computations/mixed_ranks_certificate.json}が
その有限検査を再現する。AS正則性を数値表から認定したのではなく、
AS条件から導いた有限個の関係数候補を厳密に排除した点を区別する。

\subsection{七周期同型とtraceの選択}
\begin{theorem}\label{thm:seven-period}
定義\ref{def:target}の全対象$A$には添字保持同型$A\simeq A(7)$が存在し、
最小正周期は1または7である。
\end{theorem}
\begin{proof}
命題\ref{prop:general-koszul}と定理\ref{thm:general-index}から7周期同型を得る。
$\Per(A)=\{s\in\ZZ:A\simeq A(s)\}$は、同型の平行移動・合成・逆により
$\ZZ$の部分群である。$7\in\Per(A)$なので最小正元は7の正の約数、すなわち1または7である。
\end{proof}
この定理は最小周期7の対象の存在を主張しない。
多項式増大と語長斉次二次性はAS添字7の証明で用いた仮定であり、省略していない。

\begin{proposition}[標準周期構造の選択独立性]\label{prop:trace-choice}
二つのtrace族から得た$\tau^\lambda$と$\tau^{\lambda'}$は、添字保持自己同型によって共役である。
裸の$\mathcal Q$の対象が同型なら、任意のFrobenius標準7周期構造を付した組も同型となる。
従って幾何学的同型類の集合では、標準周期構造の選択は余分な情報を加えない。
\end{proposition}
\begin{proof}
$\lambda'_i=t_i\lambda_i$と書けば、\eqref{eq:nakayama}から
$N^{\lambda'}_{ij}=(t_i/t_j)N^\lambda_{ij}$を得る。
各剰余類$i\bmod7$で初期値を選び、$s_{i+7}/s_i=t_i$を正負両方向へ解く。
$h^E_{ij}(a)=(s_i/s_j)a$は自己同型であり、
\[
h^E(7)N^\lambda(h^E)^{-1}=N^{\lambda'}
\]
が成り立つ。逆双対で$A$上の共役性を得る。
裸の同型には命題\ref{prop:linearization}を適用し、線型化した同型の双対でtraceを移送する。
\eqref{eq:nakayama}の一意性によって移送した周期構造は同型と両立し、
他のtraceの選択は上の共役で補正できる。
\end{proof}

\section{整域部分の標準型・一周期性・Ore構造}
\label{sec:domain-structure}
ここで$\ZZ$-整域とは、任意の$i\leq j\leq h$と
$0\ne a\in A_{ij}$、$0\ne b\in A_{jh}$について$ab\ne0$となることをいう。
これは$\mathcal Q$全体の定義へ加える仮定ではなく、その整域部分を取り出す条件である。

\begin{theorem}[整域部分の構造]\label{thm:domain-structure}
$A\in\mathcal Q$を$\ZZ$-整域とする。このとき次が成立する。
\begin{enumerate}
\item $A\in\mathcal S=\mathcal H$であり、各行の二変数Hilbert級数は
$((1-ut)^2(1-ut^2)(1-ut^3))^{-1}$である。
\item 語長を保つ1周期同型$A\simeq A(1)$が存在する。
従って$A\simeq\check D$となる語長二次の連結graded代数$D$が存在する。
\item $D$は、次数1の二生成元による量子平面またはJordan平面$B$を用いて
\[
D=B[z;\sigma][w;\tau,\delta],\qquad
\deg z=2,\quad\deg w=3
\]
と表せる。$B_1$上の可逆写像$\sigma,\mu$、$c\in k^*$、線型形式$\lambda$を用い、
\[
zv=\sigma(v)z,\qquad
wv=\mu(v)w+\lambda(v)z^2,\qquad wz=czw
\quad(v\in B_1)
\]
と書ける。$D$と$A$の左右noetherian性も従う。
\end{enumerate}
標準Hilbert級数・1周期性・noetherian性を追加仮定する必要はない。
\end{theorem}
\begin{proof}[証明の所在と構造]
完全証明は\path{papers/nc-p1123.tex}にある。
\texttt{idx:thm:domain-standard}では、整域性から関係数を
$(r_2,r_3,r_4,r_5)\leq(1,2,2,1)$と抑え、
七行の語長Euler式と非零jump 1生成元による行間の単射から、各上界の等号を導く。
\texttt{idx:lem:normal-family}は重み2の生成元族$z_i$の正規性と
全代数の2周期同型を与える。7周期と組み合わせた
\texttt{idx:thm:domain-period-one}から1周期性を得る。
周期の合成に互いの可換性は不要である。
1周期からgraded代数への復元は\cite[Lemma 3.4]{VdB}の構成による。

第2段のOre表示は、正則正規元$z$による商のAS次元を3へ下げ、
三生成元商のsyzygyとAS対合から$w$の正規性を導く
\texttt{idx:lem:three-generator}と\texttt{idx:thm:domain-ore}による。
正規元商の定理は重み次数2に対してではなく、$z$の\emph{語長次数1}に適用する
\cite[Lemma 1.1, Corollary 1.2]{RZ}。
商の整域性やnoetherian性を途中で仮定せず、Ore表示を得た後に導く。
\end{proof}

\begin{proposition}[重み2生成元族による整域性の判定]\label{prop:z-regular-criterion}
$A\in\mathcal Q$の語長斉次二次表示を一つ選び、
その重み2・語長1の非零生成元を$z_i\in A_{i,i+2}$とする。
次の条件は同値である。
\begin{enumerate}
\item $A$は$\ZZ$-整域である。
\item 全ての$i$で、$z_i$による左右の乗法が合成可能な全成分上で単射である。
すなわち$j\geq i+2$、$h\leq i$に対して
\[
A_{i+2,j}\longrightarrow A_{i,j},\quad a\longmapsto z_i a,
\qquad
A_{h,i}\longrightarrow A_{h,i+2},\quad b\longmapsto bz_i
\]
が単射である。
\item $i=0,\ldots,6$の七つの$z_i$について上の左右単射性が成立する。
\end{enumerate}
これは七つの生成元への還元であり、有限次数の乗法行列だけによる判定ではない。
\end{proposition}
\begin{proof}[証明の所在と要点]
完全証明は\path{papers/nc-p1123.tex}の\texttt{idx:prop:z-regular-criterion}にある。
整域なら左右単射性は明らかである。
逆向きでは、$z_i$の左右正則性だけで重み3と5の関係の各射影が単射となり、
AS対合と合わせて$(r_2,r_3,r_4,r_5)\leq(1,2,2,1)$を得る。
語長行Hilbert級数を$h_i(u)$とすると、左乗法は$h_i(u)\geq u h_{i+2}(u)$を与える。
$2$は$\ZZ/7\ZZ$を生成するので七行は互いに比較でき、Euler欠損の議論が全関係数の等号を強制する。
正規元族の構成でも必要な約分は$z_i$だけによるため、2周期から1周期へ進める。

復元したgraded代数$D$では、$z$は正則正規で語長1である。
\cite[Lemma 1.1]{RZ}とReesの補題により$D/(z)$はAS正則次元3であり、
\cite[Lemma 1.3(c)]{RZ}から$D$はnoetherianとなる。
標準語長Hilbert級数は$(1-u)^{-4}$なのでGK次元は4である。
従って\cite[Lemma 1.3]{LPWZ}を適用して$D$、次いで$A$の整域性を得る。
この段階では$D/(z)$の整域性、重み2関係の二因子階数、Ore表示を先取りしない。
7周期同型は語長を保ち、各$kz_i$を$kz_{i+7}$へ移すので、左右正則性は7周期である。
\end{proof}

\begin{corollary}[非整域部分への境界]\label{cor:domain-boundary}
$A\in\mathcal Q$について、最小正周期が7であること、
または重み4関係の$zz$射影階数$\rho_i$が$i$によって変わることのいずれかが起きれば、
少なくとも一つの$z_i$について$z_i a=0$または$bz_i=0$となる合成可能な非零元$a,b$が存在する。
整域部分では$\rho_i$は定数である。
\end{corollary}
\begin{proof}
整域なら定理\ref{thm:domain-structure}によって標準型かつ1周期である。
$\rho_i$は生成元変更で不変なので、1周期同型により$\rho_{i+1}=\rho_i$となる。
命題\ref{prop:z-regular-criterion}から$z_i$に関する結論を得る。
統合論文\texttt{idx:cor:domain-boundary}も参照。
\end{proof}

この定理は整域部分の\emph{構造の網羅性}を与える。
異なる1周期構造から復元したgraded代数は一般にgraded同型ではなく、
Zhang twistまでの同一視が必要である\cite[Theorem 1.2 = Corollary 4.4]{Sierra}。
この同一視を含む同型の存在は、次の有限条件で判定できる。
また、$\mathcal Q$全体が整域であることは証明しておらず、
零因子を持つ部分の存在自体を独立した問題として残す。

\subsection{裸の添字保持同型の有限判定}
有効な整域Ore表示を与える二つのデータ$(B,\sigma,\mu,c,L)$と$(B',\sigma',\mu',c',L')$に対し、
$B=T(V)/(kr)$、$B'=T(V')/(kr')$と書く。
ここで$L:V\to k$は最終段の導分を$D(v)=L(v)z^2$で記述する。
生成元空間への制限を同じ記号で表し、
\begin{gather*}
T=\mu\sigma^{-1},\qquad J=c^{-1}\mu^2\sigma^{-3},\\
T'=\mu'\sigma'^{-1},\qquad J'=c'^{-1}\mu'^2\sigma'^{-3},\qquad
P_i=T'^{-i}PT^i
\end{gather*}
と置く。
\begin{theorem}[同型の存在の有限条件]\label{thm:finite-isomorphism-summary}
対応する二つの付随$\ZZ$-代数が添字保持同型となることと、
可逆写像$P:V\to V'$と$\kappa\in k^\times$で
\[
\begin{gathered}
J'P=PJ,\qquad(P\otimes P_1)r\in k^\times r',\\
L'P_3=\kappa L,\qquad L'P_4=\kappa(c/c')L
\end{gathered}
\]
を満たすものが存在することは同値である。
\end{theorem}
完全証明は\path{papers/nc-p1123.tex}の\texttt{idx:thm:finite-isomorphism}にある。
語長の線型化、頂点スカラーの正規化、二つのshiftから得る
$P_i=T'^{-i}PT^i$の漸化式が、全ての添字の条件を上の有限系へ帰着する。
これは非線型同型も許した\emph{同型の存在}の必要十分条件であり、
全自己同型群や一般の軌道の標準形、moduli stackを記述するものではない。
個別結果の文献上の新規性の最終確認は未了である。

\section{関係の七窓による有限表示}
まず標準型$\mathcal S$の場合を記す。
$V_i^{(a)}=V_{i,i+a}$とし、$a=1,2,3$、$i\bmod7$で生成元空間を固定する。
\[
W_{i,d}=\bigoplus_{a+b=d}V_i^{(a)}\otimes V_{i+a}^{(b)}
\]
は、語長2・重み$d$の全経路空間である。
\begin{center}
\begin{tabular}{c|rrrrr}
\toprule
重み$d$ & 2 & 3 & 4 & 5 & 6\\
\midrule
$\dim W_{i,d}$ & 4 & 4 & 5 & 2 & 1\\
$\dim R_{i,i+d}$ & 1 & 2 & 2 & 1 & 0\\
関係Grassmannianの次元 & 3 & 4 & 6 & 1 & 0\\
\bottomrule
\end{tabular}
\end{center}
従って七窓の表示空間と線型生成元変更群は
\begin{align*}
\mathfrak P_7&=\prod_{i=0}^{6}
\left(\Gr(1,4)\times\Gr(2,4)\times\Gr(2,5)\times\Gr(1,2)\right),\\
G_7&=\prod_{i=0}^{6}\left(\GL_2\times\GG_m^2\right),
\qquad \dim\mathfrak P_7=98,\quad\dim G_7=42.
\end{align*}
各因子での作用は、$W_{i,d}$の経路ごとに
$g_i^{(a)}\otimes g_{i+a}^{(b)}$を作用させるものである。

\begin{theorem}[標準型の有限表示]\label{thm:finite-window}
$\mathfrak X_7$を、$\mathfrak P_7(k)$のうち生成される$\ZZ$-代数が標準型であり、
厳密な7移動があるtrace族から来る標準周期構造となる点の集合とする。
標準型の裸の幾何学的同型類の集合は$\mathfrak X_7/G_7(k)$と全単射する。
これは点集合についての主張であり、相対族の下降やmoduli stackの同定を含まない。
\end{theorem}
\begin{proof}
標準周期同型を選び、始点$0,\ldots,6$で各jumpの生成元の基底を選ぶ。
周期同型で全整数へ基底を延長すると関係データは厳密7周期となり、$\mathfrak X_7$の点を得る。
逆に七窓の関係を整数頂点列上へ周期延長してテンソル代数を割れば対象を復元する。
標準周期同型と可換する語長保存同型は七つの基底変更で決まる。
裸の同型は命題\ref{prop:linearization}と命題\ref{prop:trace-choice}によって
そのような同型へ補正できる。従って軌道は裸の同型類と一致する。
\end{proof}

定理\ref{thm:all-standard-summary}により$\mathcal Q=\mathcal S$であり、
全対象が$r_2=r_5=1$、$r_3=r_4=2$、$r_6=0$という同じ階数表を持つ。
従って上の$\mathfrak X_7/G_7(k)$は全$\mathcal Q$の幾何学的同型類を与え、
次元98の表示空間が全対象の表示を収容する。
この空間内のAS/Koszul locusの必要十分条件、軌道の標準形、
全モジュライの次元や成分を決めたとはしない。

標準周期条件自体は有限対合行列で書ける。
厳密7周期表示のもとで七つのtop空間のtraceを選び、厳密7移動によって全整数へ延長する。
逆に標準周期が厳密移動と一致すれば、\eqref{eq:nakayama}をtop要素と恒等射に適用して
$\lambda_{i+7}(N(a))=\lambda_i(a)$を得るため、この選び方で十分である。
$u_\alpha\in E_{i,i+a}^1$、$v_\beta\in E_{i+a,i+7}^3$を基底とし、
$u'_\gamma$を厳密7移動した基底とする。
\[
P_{i,a}=\bigl(\lambda_i(u_\alpha v_\beta)\bigr),\qquad
Q_{i,a}=\bigl(\lambda_{i+a}(v_\beta u'_\gamma)\bigr).
\]
式\eqref{eq:nakayama}は$P_{i,a}^{\mathsf T}=Q_{i,a}N_{i,a}$となるため、
標準周期条件は$P_{i,a}^{\mathsf T}=Q_{i,a}$である。
\textbf{jump $a=1,2,3$を全て検査する。}
jump 1だけの検査では重み2と3の生成元の移動を制御できない。

\begin{remark}\label{rem:finite-not-koszul}
語長$0,\ldots,5$における二次双対の次元表と相補積は有限線型代数で検査できる。
語長5では重み15までを調べればよく、語長5の全ての積が零なら高い語長も零となる。
しかしこの有限Frobenius検査だけでは元代数のKoszul複体の正確性は従わない。
本稿は$\mathfrak X_7$の有限な必要十分方程式を完成させたとは主張しない。
また$98-42=56$をモジュライ次元とはしない。
\end{remark}

\section{対角量子族と重み付きtwist}
$w=(1,1,2,3)$、$q_{rs}\in k^*$ ($0\leq r<s\leq3$) として
\[
B(q)=k\langle x_0,x_1,x_2,x_3\rangle/
(x_sx_r-q_{rs}x_rx_s)_{r<s}
\]
に重み$\deg x_r=w_r$を与える。
これはAKO\cite[Proposition 2.6]{AKO}の既知の正則族である。
全ての$k$-点で順序単項式$x_0^{n_0}\cdots x_3^{n_3}$を基底に持つ。
三つの異なる添字$k>j>i$の重なりは、どちらの順で簡約しても
$q_{ij}q_{ik}q_{jk}x_ix_jx_k$となる。
二次双対は量子外積代数で、相補単項式の積係数は非零なのでFrobeniusである。
従って$\check B(q)$は標準型に属する。

\begin{proposition}[対角twistの正規化]\label{prop:twist-normalization}
対角自己同型$\sigma(x_r)=\lambda_r x_r$による重み付きZhang twistを
$a*b=a\sigma^{\deg a}(b)$で定義すると
\[
q'_{rs}=q_{rs}\frac{\lambda_r^{w_s}}{\lambda_s^{w_r}}.
\]
この作用の軌道は、次の三つのLaurent単項式で完全に記述される。
\begin{equation}\label{eq:three-invariants}
r_{12}=\frac{q_{01}^2q_{12}}{q_{02}},\qquad
r_{13}=\frac{q_{01}^3q_{13}}{q_{03}},\qquad
r_{23}=\frac{q_{02}^3q_{23}}{q_{03}^2}.
\end{equation}
特に各対角族の$\ZZ$-代数は、$q_{01}=q_{02}=q_{03}=1$で、
残りの係数が$(r_{12},r_{13},r_{23})$であるものと同型である。
\end{proposition}
\begin{proof}
$x_s*x_r=\lambda_r^{w_s}q_{rs}x_rx_s$と
$x_r*x_s=\lambda_s^{w_r}x_rx_s$を比較すると変換式を得る。
$\lambda_0=1$, $\lambda_j=q_{0j}$ ($j=1,2,3$) と選べば最初の三係数は1となり、
残りは\eqref{eq:three-invariants}となる。
変換式への代入でこれらが不変であることが分かり、同じ三係数を持つ二点は
同じ正規化へ移るので軌道の完全不変量である。

さらに$a\in B(q)_{j-i}$に対して
$\phi_{ij}(a)=\sigma^i(a)$と置けば、
$\sigma^i(a\sigma^{j-i}(b))=\sigma^i(a)\sigma^j(b)$より
$\check{B(q)^\sigma}\simeq\check B(q)$を得る。
従って正規化は添字保持同型も与える。
\end{proof}

\begin{theorem}[対角族の全特殊値における裸の同型]\label{thm:diagonal-indexed-summary}
正規化した三係数を$(a,b,c)=(r_{12},r_{13},r_{23})$と書く。
二つの三係数が同じ添字保持同型類を与えることと、
同一であるか、対合
\[
s(a,b,c)=(a^{-1},b^{-1},ca^3/b^2)
\]
で移り合うことは同値である。
従ってこの族の裸の同型類の集合は$(\GG_m)^3(k)/\langle s\rangle$であり、
そのaffine有限商は既約かつ次元3である。
固定点集合は$a=1$, $b=\pm1$, $c\in k^\times$である。
\end{theorem}
完全証明は統合論文\texttt{diag:thm:indexed-isomorphism}、
\texttt{diag:cor:quotient}にある。一般線型変更と固有値の一致する例外値も含む。
命題\ref{prop:twist-normalization}の三係数を求めた後、
上の二つの場合を比較すれば、元の六係数の族の裸の同型が判定できる。
これは対角族の商の記述であり、全ASモジュライの成分やstackの同定ではない。
個別結果の文献上の新規性は未確認である。

\section{非対角Ore族と最初の不変分岐}
\subsection{全パラメータでの正則性}
\begin{theorem}\label{thm:ore}
$a,b\in k$として
\[
A_{a,b}=k\langle x,y,z,t\rangle/
\left(yx-xy,\ zx-xz,\ zy-yz,\ tx-xt-az^2,\ ty-yt-bz^2,\ tz-zt\right)
\]
に重み$(1,1,2,3)$を与える。この族の全てのファイバーは語長Koszulで、
四次元AS正則、AS添字7であり、$\check A_{a,b}$は標準型に属する。
族内の添字保持同型類は$\check A_{0,0}$と$\check A_{1,0}$のちょうど二つである。
後者は、どの対角量子多項式$\ZZ$-代数とも同型でない。
\end{theorem}
\begin{proof}[正則性の証明]
$k[x,y,z]$上の導分$\delta(x)=az^2$, $\delta(y)=bz^2$, $\delta(z)=0$を使うと
$A_{a,b}=k[x,y,z][t;\delta]$である。
順序$t>z>y>x$のdegree-lex順で六関係を簡約規則として向けると、全ての右辺の項は左辺より小さい。
包含曖昧性はなく、重なりは次の四つで尽くされる。
\begin{center}
\begin{tabular}{c|l}
\toprule
語 & 左右の共通の正規形\\
\midrule
$zyx$ & $xyz$\\
$tyx$ & $xyt+ayz^2+bxz^2$\\
$tzx$ & $xzt+az^3$\\
$tzy$ & $yzt+bz^3$\\
\bottomrule
\end{tabular}
\end{center}
Diamond lemma\cite[Theorem 1.2]{Bergman}により$x^\alpha y^\beta z^\gamma t^\delta$が基底となる。
二次Gr\"obner基底を持つので、語長についてKoszulである\cite[Theorem 5.3]{Priddy}。

二次双対$E=A_{a,b}^!$を直接計算する。
生成元を$X,Y,Z,T$とすると、異なる生成元は反可換であり、
\begin{equation}\label{eq:ore-dual}
X^2=Y^2=T^2=0,\qquad Z^2=aXT+bYT
\end{equation}
である。例えば$Z^2+aTX+bTY$は六つの原関係に直交し、
反可換関係と合わせて\eqref{eq:ore-dual}を得る。
全16次元の二次テンソル空間で、原関係は6次元なので、これら10本の独立関係で双対を尽くす。

$C=\Lambda(X,Y,T)$とし、$\sigma$をparity自己同型とする。
$c=(aX+bY)T$は$C$の偶中心元で$\sigma(c)=c$だから
\[
E\simeq C[Z;\sigma]/(Z^2-c).
\]
$Z^2-c$は中心的なmonic多項式なので割算により$C$上で$1,Z$が基底となる。
従って
$X^\epsilon Y^\eta Z^\theta T^\kappa$ ($\epsilon,\eta,\theta,\kappa\in\{0,1\}$)
が$E$の基底で、語長次元は$(1,4,6,4,1)$、topの重みは7である。

$XYZT$の係数を取るtraceを考える。
二つの基底単項式の積に$Z$が二度現れれば、\eqref{eq:ore-dual}で$Z$が消え、top係数は零となる。
それ以外の積は外積代数と同じ相補単項式対合を与える。
よって各語長$p$と$4-p$の積対合は完全であり、重みも$d$と$7-d$を対にする。
従って$E$はFrobeniusである。

Koszul複体は正確で、その項は$E^*$から与えられる。
このFrobenius対合でKoszul複体の双対を反転した複体と同一視すると、
左右のAS双対性を得る。Koszul双加群解消も長さ4なのでglobal dimensionは4以下であり、
$E^4\ne0$よりちょうど4となる。
topの重みからAS添字は7である。基底単項式により成長は多項式的であり、
各シフトは\eqref{eq:resolution}と一致する。
connected graded代数のKoszul--Frobenius判定\cite[Corollary D]{LPWZKoszul}の具体例でもある。
同文献のAS定義に含まれない多項式成長は、ここではPBW基底から別途確認した。
同型についての残りの主張は次の命題から従う。
\end{proof}

\subsection{関係の射影階数}
\begin{proposition}\label{prop:rho}
語長斉次二次表示を持つ対象について、
\[
\rho_i(A)=\rk\left(R_{i,i+4}\longrightarrow
V_{i,i+2}\otimes V_{i+2,i+4}\right)
\]
は添字保持同型不変である。右辺の射は二次経路空間の直和分解への射影である。
特に$\rho_i\in\{0,1\}$であり、$\mathcal Q$の全対象で$\rho_{i+7}=\rho_i$である。
\end{proposition}
\begin{proof}
語長付随次数において$R=\ker(V\otimes_{\ZZ}V\to I^2/I^3)$は内在的である。
添字を保つ生成元変更は中間頂点$i+1,i+2,i+3$ごとの経路直和を保つので、射影階数も保つ。
命題\ref{prop:linearization}により、これは非線型同型を許しても変わらない。
値域の次元は1である。$\mathcal Q$の対象では定理\ref{thm:seven-period}を適用する。
\end{proof}

\begin{proof}[定理\ref{thm:ore}の同型分類の証明]
$(a,b)\ne(0,0)$なら、$kx+ky$の基底を取り直して
$\delta(x')=z^2$, $\delta(y')=0$とできる。
従って$A_{a,b}\simeq A_{1,0}$は重みを保つgraded同型である。
一方、$\check A_{a,b}$に対する射影の行列は$(-a,-b)$なので、
零点では$\rho_i=0$、非零点では$\rho_i=1$である。
よって両者はindexed同型でも異なり、族内にちょうど二つの同型類がある。
対角量子族の重み4関係は$x_rt$と$tx_r$ ($r=0,1$) だけを含み、
$zz$成分がないため常に$\rho_i=0$である。
従って非零Ore族はそのどの対象とも同型でない。
\end{proof}

\begin{remark}
この二同型類の分類は指定した二パラメータ族に関する完全な結果であり、
$\rho_i=1$の全標準対象の分類ではない。
零のlocusにも対角族以外の対象があるかは、全軌道を調べる問題として残る。
非定数の$\rho$パターンは系\ref{cor:domain-boundary}により非整域部分にしか現れず、
そのような対象の存在は未解決である。
またこの族の構成は既知のOre手法の範囲にあり、新規性は未確認である。
\end{remark}

\section{対角skew Oreチャートの表示分類}
既知の二パラメータOre族を含む、対角量子三生成元部分の上の
一般の斉次skew Oreチャートを解析した。
その必要十分条件は四つの係数式であり、表示パラメータ空間の台集合は
次元$6,5,5,4$の四つの既約成分に分かれる。
全ファイバーの正則性と関係の$zz$射影階数による非同型性を含め、
完全な証明は統合公表用論文\path{papers/nc-p1123.tex}に置く。
これらは表示空間の成分であり、全モジュライの成分とは主張しない。
\section{次数付き部分の段階的Ore分類}
この節では基礎体を$\mathbb C$とし、重みを保つgraded同型を用いる。
以下の定理の対象は、語長二次性を仮定しないgraded代数である。
語長二次表示を持つものの付随$\ZZ$-代数は$\mathcal H$に属する。
逆に$\mathcal Q$の整域部分は、定理\ref{thm:domain-structure}により
この節の語長二次Ore表示で全て実現される。

\begin{theorem}[追加仮定付きgraded分類]\label{thm:graded-ore-summary}
$A$を連結・局所有限・左右noetherianな四次元AS正則graded $\mathbb C$-代数とし、
最小生成元次数を$(1,1,2,3)$、Hilbert級数を
\[
H_A(t)=\frac{1}{(1-t)^2(1-t^2)(1-t^3)}
\]
とする。このとき、次数1の二生成元$x,y$からなる量子平面またはJordan平面$B$と、
次数2、3の最小生成元$z,w$を選び、
\[
A=B[z;\sigma,\delta][w;\tau,D]
\]
と表せる。ここで$\sigma$は$B$のgraded自己同型、$\delta$は次数2の$\sigma$-導分、
$\tau$は$C=B[z;\sigma,\delta]$のgraded自己同型、$D$は次数3の$\tau$-導分である。
逆に、これらのデータから作る全ての代数が上の仮定を満たす。
\end{theorem}
証明は\path{papers/nc-p1123.tex}の
\texttt{thm:ore-exhaustion}にある。
最小解消の双対性から関係の両側の係数階数を求め、
指定Hilbert級数を用いて順序単項式の独立性を示す。
正則性の逆向きには\cite[Theorem~3.3]{LWW}を適用し、成長条件をPBW基底で確認する。

同稿の\texttt{thm:coefficient-data}は許容データを有限の係数条件で記述し、
\texttt{thm:graded-isomorphisms}は全graded同型を
\[
u\longmapsto\varphi(u)\quad(u=x,y),\qquad
z\longmapsto\alpha z'+a_2,\qquad
w\longmapsto\beta w'+b_3
\]
とその必要十分条件で与える。
ここで$\varphi:B\simeq B'$、$\alpha,\beta\in\mathbb C^\times$、
$a_2\in B'_2$、$b_3\in C'_3$である。
中心的な次数3最小生成元を選べる場合は$A=C[w]$となる
（同稿\texttt{cor:central-generator}）。

語長二次表示に適合する生成元を選ぶと、
$\delta=0$、$\tau|_B=\mu$、$\tau(z)=cz$、
$D(x),D(y)\in\mathbb C z^2$、$D(z)=0$となる。
この部分の全必要十分条件とgraded同型の存在判定は、同稿の
\texttt{cor:quadratic-chart}、\texttt{prop:quadratic-isomorphisms}にある。
対角性を仮定しないため、可換平面・$q=-1$の量子平面・Jordan平面の
非対角自己同型も含む。量子平面の上で両自己同型を対角に制限すると、前節の四係数式を回収する。

定理\ref{thm:graded-ore-summary}のnoetherian性とHilbert条件は
$\mathcal Q$全体に加えた仮定ではなく、この節の適用範囲である。
また裸の添字保持同型ではZhang twistも同一視されるため、
上のgraded同型条件だけから全indexed同型の商は得られない。
語長二次の整域部分には定理\ref{thm:finite-isomorphism-summary}の有限な必要十分条件がある。

次数1の最小生成元を$n$個に増やした場合の構造定理は
\path{papers/nc-p1123.tex}にある。
左右noetherian・整域・AS次元$n+2$・所期の重み付き多項式Hilbert級数・
語長二次性の下では、任意の$n$について
$A=B[z;\sigma][w;\tau,D]$という表示が網羅的である。
$B=k\langle A_1\rangle$はAS正則次元$n$、Hilbert級数$(1-t)^{-n}$を持つ。
同稿は正規元の商の二変数Hilbert級数とnoetherian性を使って任意次元の網羅性を証明し、
一般の斉次Ore部分族の係数条件と全graded同型も記述する。
語長二次性を外した全$n$の網羅性は未証明であり、
最小生成元次数だけを指定した場合には大域次元が無界となる明示族がある。

\section{残る課題と検証の範囲}
整域部分については、標準型の強制、1周期性、量子平面またはJordan平面からの
2段Ore表示が得られ、Hilbert級数とnoetherian性は追加仮定でなくなった。
次の構造上の核心は、各$z_i$の左右乗法の単射性をAS条件から導けるかである。
命題\ref{prop:z-regular-criterion}により、この問題は$\mathcal Q$全体の整域性と同値である。
七つの代表$z_0,\ldots,z_6$に対して最小語長の左・右零化元を仮定し、
最小解消、AS対合、関係の射影とそのsyzygyから矛盾を導けるかを調べる。
有限次数で零化元を検出しなかったことだけでは正則性を認定しない。
最小周期7または非定数の$\rho$は$z_i$に零因子がある部分に限られる。
七窓の関係射影が退化する枝と、高い語長に初めて零化元が現れる可能性をともに扱う。
有限双対のFrobenius性とKoszul複体の正確性は独立に扱う。

数値型については、定理\ref{thm:outer-relation-ranks}で$r_2=r_5=1$と
純jump 1部分のHilbert係数$n+1$を得た。
混合関係数も定理\ref{thm:all-standard-summary}により$r_3=r_4=2$に確定した。
一本の重み2関係tensorの行列階数2や、$B$の整域性・AS正則次元2もまだ結論していない。
研究ノートの非負Euler係数表は、実現する二次代数の追加構造を使う必要を示す障害例である。

分類上は、整域部分の裸の同型の存在に関する必要十分条件は
定理\ref{thm:finite-isomorphism-summary}で有限系に帰着した。
対角族の軌道は定理\ref{thm:diagonal-indexed-summary}で全特殊値を含めて記述した。
残るのは一般Oreデータの軌道の標準形、成分・特殊化・安定化群とモジュライ幾何である。
対角族の明示的な有限商と非対角基準族を回帰例として、有限判定式の解集合を調べる。
degree 1 del Pezzoの反標準環や$E_8$量子化との比較では、
四つの最小生成元と語長二次性が保たれるかを先に確かめる。
関連文献の確認範囲とこの条件の障害は研究ノートへ置く。

\path{computations/weighted_baseline.py}は、二変数Euler式、重み24までのHilbert係数、
関係空間の次元、対角twistの整数指数行列と正規化、各族の四つの三重重なりを形式的に検算する。
構造定理の解析的な証明は有限次数のHilbert照合を根拠としない。
一般対象の数値標準性の最終段階には、別の\path{computations/verify_mixed_ranks.py}による
有限整数候補の網羅的な行列式計算を用いる。
証明の記載・独立監査・計算・組版・保存の状態は\path{research/validation.json}で区別する。
全分類、形式化、第三者査読、新規性の最終確認は完了していない。

\begin{thebibliography}{99}
\bibitem{AKO}
D.~Auroux, L.~Katzarkov and D.~Orlov,
\emph{Mirror symmetry for weighted projective planes and their noncommutative deformations},
Annals of Mathematics \textbf{167} (2008), 867--943.
\href{https://doi.org/10.4007/annals.2008.167.867}{doi:10.4007/annals.2008.167.867};
\href{https://arxiv.org/abs/math/0404281}{arXiv:math/0404281}.
\bibitem{Sierra}
S.~J.~Sierra, \emph{G-algebras, twistings, and equivalences of graded categories},
Algebras and Representation Theory \textbf{14} (2011), 377--390.
\href{https://doi.org/10.1007/s10468-009-9193-y}{doi:10.1007/s10468-009-9193-y};
\href{https://arxiv.org/abs/math/0608791}{arXiv:math/0608791}.
\bibitem{VdB}
M.~Van den Bergh, \emph{Non-commutative quadrics},
International Mathematics Research Notices \textbf{2011}, no.~17, 3983--4026.
\href{https://doi.org/10.1093/imrn/rnq234}{doi:10.1093/imrn/rnq234};
\href{https://arxiv.org/abs/0807.3753}{arXiv:0807.3753}.
\bibitem{MVS}
R.~Mart\'inez-Villa and \O.~Solberg,
\emph{Artin--Schelter regular algebras and categories},
Journal of Pure and Applied Algebra \textbf{215} (2011), 546--565.
\href{https://doi.org/10.1016/j.jpaa.2010.06.006}{doi:10.1016/j.jpaa.2010.06.006};
\href{https://arxiv.org/abs/0902.1119}{arXiv:0902.1119}.
\bibitem{LWW}
L.~Liu, S.~Wang and Q.~Wu, \emph{Twisted Calabi--Yau property of Ore extensions},
Journal of Noncommutative Geometry \textbf{8} (2014), 587--609.
\href{https://doi.org/10.4171/JNCG/165}{doi:10.4171/JNCG/165};
\href{https://arxiv.org/abs/1205.0893}{arXiv:1205.0893}.
\bibitem{LPWZ}
D.-M.~Lu, J.~H.~Palmieri, Q.-S.~Wu and J.~J.~Zhang,
\emph{Regular algebras of dimension 4 and their $A_\infty$-Ext-algebras},
Duke Mathematical Journal \textbf{137} (2007), 537--584.
\href{https://doi.org/10.1215/S0012-7094-07-13734-7}{doi:10.1215/S0012-7094-07-13734-7};
\href{https://arxiv.org/abs/math/0411497}{arXiv:math/0411497}.
\bibitem{Bergman}
G.~M.~Bergman, \emph{The diamond lemma for ring theory},
Advances in Mathematics \textbf{29} (1978), 178--218.
\href{https://doi.org/10.1016/0001-8708(78)90010-5}{doi:10.1016/0001-8708(78)90010-5}.
\url{https://math.berkeley.edu/~gbergman/papers/updates/diamond.html}.
\bibitem{Priddy}
S.~B.~Priddy, \emph{Koszul resolutions},
Transactions of the American Mathematical Society \textbf{152} (1970), 39--60.
\url{https://math.mit.edu/~hrm/palestine/priddy-koszul-resolutions.pdf}.
\bibitem{RZ}
D.~Rogalski and J.~J.~Zhang,
\emph{Regular algebras of dimension 4 with 3 generators}.
\href{https://arxiv.org/abs/1101.1998}{arXiv:1101.1998v2}.
\bibitem{LPWZKoszul}
D.-M.~Lu, J.~H.~Palmieri, Q.-S.~Wu and J.~J.~Zhang,
\emph{Koszul equivalences in $A_\infty$-algebras},
New York Journal of Mathematics \textbf{14} (2008), 325--378.
\url{https://nyjm.albany.edu/j/2008/14-16.pdf}.
\end{thebibliography}
\end{document}
